% GitHub Repo and Documentation: https://github.com/celiobjunior/resume-template
% Copyright © 2025 Celio B Junior. All rights reserved.
% 
% Licensed under the Apache License, Version 2.0 (the "License");
% you may not use this file except in compliance with the License.
% You may obtain a copy of the License at
%
%     http://www.apache.org/licenses/LICENSE-2.0
%

\documentclass[a4paper,10pt]{article}

% Layout
\usepackage{geometry}
\geometry{top=1cm,bottom=1cm,left=1cm,right=1cm}
\usepackage{parskip}
\pagestyle{empty}

% Fonts & language (XeLaTeX / LuaLaTeX only)
\usepackage{fontspec}
\usepackage{polyglossia}
\setdefaultlanguage{portuguese}
\setotherlanguage{english}

\setmainfont{TeX Gyre Termes}
\setsansfont{TeX Gyre Heros}
\setmonofont{TeX Gyre Cursor}

% Colors & links
\usepackage{xcolor}
\definecolor{muted}{gray}{0.20}

\usepackage{hyperref}
\hypersetup{
  pdftitle={CV Marlon Correia Bomfim Junior},
  pdfauthor={Marlon Correia Bomfim Junior},
  colorlinks=true,
  linkcolor=black,
  urlcolor=black
}

% Section formatting
\usepackage{titlesec}
\setcounter{secnumdepth}{0}

\titleformat{\section}
  {\sffamily\Large\bfseries}
  {}
  {0pt}
  {}

\titleformat{\subsection}
  {\sffamily\bfseries}
  {}
  {0pt}
  {}

% Helpers
\newcommand{\sectionend}{
  \par\addvspace{0.75em}
  \hrule
  \addvspace{0.75em}
}

\newcommand{\company}[1]{{\bfseries #1}}
\newcommand{\place}[1]{{\itshape #1}}
\newcommand{\role}[1]{{\itshape\color{muted}#1}}
\newcommand{\daterange}[1]{{\small\bfseries\color{muted}#1}}
\newcommand{\tech}[1]{{\ttfamily #1}}

\setlength{\emergencystretch}{2em}
\begin{document}

% --- CABEÇALHO ---
\begin{flushleft}
    {\LARGE \textbf{Marlon Correia Bomfim Junior}} 
    \\ [0.1cm]
    Praia Grande, SP, Brasil
    \quad | \quad
    {\ttfamily \href{https://github.com/marlaoza}{Github}
    \quad | \quad
    \href{https://www.linkedin.com/in/marlon-correia-bomfim-junior-36a1a9230}{Linkedin} 
    \quad | \quad
    \href{mailto:marlonbomfim.jr@gmail.com}{marlonbomfim.jr@gmail.com} 
    }
\end{flushleft}

% --- SEÇÕES ---

\section{Experiência}
    \subsection*{\texorpdfstring{
      \company{OpenQuant}
      \hfill
      \place{Praia Grande, SP, Brasil}
    }{
      OpenQuant -- Praia Grande, SP, Brasil
    }}
    \par
    \role{Desenvolvedor Júnior}
    \hfill
    \daterange{Fev 2024 -- Atual}
        
    Atuação Full Stack no desenvolvimento da plataforma OpenQuant.  
    Criação de componentes e interfaces responsivas e modernas em React.  
    Desenvolvimento e planejamento de APIs escaláveis em .NET, com migrations e paradigmas de injeção de dependências.  
    Manejo de dados em PostgreSQL e hot cache com Redis.  
    Dockerização e deploy de ambientes integrados com Kubernetes.


    \subsection*{\texorpdfstring{
      \company{MCB Solution}
      \hfill
      \place{Remoto}
    }{
      MCB Solution -- Remoto
    }}
    \par
    \role{Desenvolvedor Júnior}
    \hfill
    \daterange{Fev 2023 -- Fev 2024}
        
   Desenvolvimento de diversas plataformas de automação para clientes em Angular e React. Montagem de fluxos e ferramentas de automação em Python e C\#. Integração de robôs atuando em tarefas reais.

    \subsection*{\texorpdfstring{
      \company{MCB Solution}
      \hfill
      \place{Remoto}
    }{
      MCB Solution -- Remoto
    }}
    \par
    \role{Estagiário de Desenvolvimento}
    \hfill
    \daterange{Ago 2022 -- Fev 2023}
        
    Desenvolvimento de algoritmos de automação para geração de relatórios, extratos e notas integradas a ferramentas. Planejamento de arquiteturas de ASP.NET MVC flexíveis gerenciando ambientes de BackOffice. Desenvolvimento de interfaces WPF+ para robôs de boletagem.


  

\sectionend

% ------

\section{Educação}

    \subsection*{\texorpdfstring{
      \company{Universidade Santa Cecília}
      \hfill
      \place{Santos, SP, Brasil}
    }{
      Universidade Santa Cecília -- Santos, SP, Brasil
    }}
    \par
    \role{Bacharelado em Engenharia da Computação}
    \hfill
    \daterange{Fev 2023 - Dez 2027}

    \sectionend

\section{Projetos}

    \subsection*{\texorpdfstring{
      \company{Oficina Curto Circuito}
      \hfill
      \place{Guarujá, SP, Brasil}
    }{
      Oficina Curto Circuito -- Guarujá, SP, Brasil
    }}
    \par
    \role{Professor de Eletrônica}
    \hfill
    \daterange{Dez 2025 - Atual}
        
    Projeto em parceria com a Escola Estadual Vicente de Carvalho, onde ministro, juntamente com meus colegas, 
    uma oficina de eletrônica e automação para o oitavo e nono ano, ensinando desde circuitos básicos com leds 
    até projetos mais robustos com microcontroladores, a fim de compartilhar conhecimento e capacitá-los a montar seus próprios projetos.

    \subsection*{\texorpdfstring{
    \company{Salve a Nave - Plataforma Educativa}
    \hfill
    \place{Santos, SP, Brasil}
    }{
      Salve a Nave - Plataforma Educativa -- Santos, SP, Brasil
    }}
    \par
    \role{Desenvolvedor de aplicação C\#}
    \hfill
    \daterange{Ago 2024 - Atual}
        
    FINALISTA COBRIC 2025 - Software interativo que apresenta uma série de desafios em formato de minijogos que se encaixam no contexto lúdico de uma nave espacial, feito com C\#, como intervenção educacional, criando um espaço para crianças desenvolverem habilidades 
    motoras, sociais e de raciocínio lógico, além de aprimorarem o conhecimento passado em sala de aula. 
    Conta com uma ferramenta para criação de fases, onde professores e/ou responsáveis tem controle das matérias abordadas.
    
    \textbf{Links}
    
    {\ttfamily
    \par
    \begin{tabular*}{\linewidth}{@{\extracolsep{\fill}} c c c }
    \href{https://zenodo.org/records/17778281}{DOI} & 
    \href{https://www.conic-semesp.org.br/anais/files/2025/trabalho-1000014248.pdf?1770686772}{CONIC 2025} &
    \href{https://ojs.unisanta.br/COB/article/view/3203/3079}{COBRIC 2025}
    \end{tabular*}
    }


    \subsection*{\texorpdfstring{
    \company{ELIB - Interprete de LIBRAS}
    \hfill
    \place{Santos, SP, Brasil}
    }{
      ELIB - Interprete de LIBRAS -- Santos, SP, Brasil
    }}
    \par
    \role{Desenvolvedor de aplicação mobile Kotlin + Swift}
    \hfill
    \daterange{Jan 2024 - Dez 2024}
        
    FINALISTA COBRIC 2024 - Aplicativo mobile nativo, desenvolvido para Android com Kotlin e iOS com Swift, que integra via TensorFlow um modelo de visão computacional criado do zero para classificação de gestos de LIBRAS.  
    O aplicativo apresenta uma interface de fácil acesso, com tradução bidirecional: uma captando gestos pela câmera e convertendo em texto, e outra recebendo input por voz ou texto e traduzindo para sinais de LIBRAS.
    
    \textbf{Links}
    
    {\ttfamily
    \par
    \begin{tabular*}{\linewidth}{@{\extracolsep{\fill}} c c c }
    \href{https://zenodo.org/records/14040753}{DOI} & 
    \href{https://www.conic-semesp.org.br/anais/files/2024/trabalho-1000011871.pdf?1770686972}{CONIC 2024} &
    \href{https://periodicos.unisanta.br/COB/article/view/2292/2315}{COBRIC 2024}
    \end{tabular*}
    }
    
\sectionend

% ------

\section{Atividades} 

    \subsection*{\texorpdfstring{
    \company{Game Jams}
    \hfill
    \place{Remoto}
    }{
      Game Jams -- Remoto
    }}
    \par
    \role{Ilustrador / Desenvolvedor}
    \hfill
    \daterange{Jan 2024 e Ago 2025}
        
    Distribuição e manejamento de tarefas para desenvolvimento completo de jogos em 7 dias. Planejamento e execução de sistemas completos, escaláveis e práticos em C\# com Unity.
    Desenvolvimento de artes do zero com ilustrações próprias.
    
    \textbf{Links}
    
    \par
    {\ttfamily
    \begin{tabular*}{\linewidth}{@{\extracolsep{\fill}} c c c }
    \href{https://arthurpalladino.itch.io/nothing-can-go-wrong-at-the-theater}{GAME JAM 01/2024} & 
    \href{https://arthurpalladino.itch.io/take-a-risk-or-get-cut-off}{GAME JAM 08/2025}
    \end{tabular*}
    }
    
\sectionend
% ------

\section{Habilidades}
    \subsection*{Desenvolvimento}
        \subsubsection*{Linguagens}
        {\ttfamily 
            Javascript (Typescript) \,|\, 
            C \,|\, 
            C++ \,|\, 
            C\# \,|\, 
            Python \,|\, 
            AVR Assembly
        }
        \subsubsection*{Embarcados e Hardware}
        {\ttfamily
        ESP32 \,|\, 
        Microcontroladores ATMega
        \subsubsection*{Frameworks}
        React \,|\, 
        Angular \,|\, 
        .NET \,|\, 
        Kotlin
        }
    \vspace{0.2em}
    \subsection*{Softwares de Criação}
    {\ttfamily
    Photoshop (Manipulação de imagens, ilustração digital) \,|\, 
    DaVinci Resolve (Edição de vídeo e efeitos especiais) \,|\, 
    Unity (Criação de jogos e ferramentas) \,|\, 
    KiCad (Prototipação de placas de circuito impresso) \,|\, 
    Fusion 360 (Prototipação de peças 3D)
    }
    \vspace{0.2em}
    \subsection*{Idiomas}
    {\ttfamily
    Português (Nativo) \,|\, 
    Inglês (Avançado) \,|\, 
    Espanhol (Básico)
    }
\sectionend
\end{document}

